% 本文に出す数値。本文ではここのマクロを使い、解析をやり直したらここだけを更新する。
% 各 block の上に出典を書く: source 名、run id、解析の path、値を作った script。
% 丸め方、区間の定義、単位の約束もここに書く。
%
% 例:
% % source: pipe_flow, runs/roughness-scan/_analysis/pressure/out/summary.csv
% %         (R20261015-0003 .. R20261015-0008), figures/pressure-loss/make.py
% \newcommand{\lossRatio}{3}   % 粗さ 2% の管と滑らかな管の圧力損失の比

% まだ得ていない値は、予想を \pred{} で書く。結果が出たら \pred を外して出典を書く。
% % 予想: 粗さに比例するとした見積もり（BRIEF C2）。roughness-scan の再計算で置き換える
% \newcommand{\stressRatio}{\pred{2.8}}

% 予想も立てられない値に置く。
\newcommand{\pending}{\textcolor{red}{[pending]}}
